
\documentclass[12pt]{article}
\usepackage[a4paper,margin=1in]{geometry}
\usepackage{amsmath,amssymb,mathtools}
\usepackage{bm}
\usepackage{siunitx}
\usepackage{microtype}
\usepackage[hidelinks]{hyperref}
\usepackage[capitalise,noabbrev]{cleveref}
\numberwithin{equation}{section}
\sisetup{separate-uncertainty=true}

% --- Macros (avoid redefining \d which is an accent in LaTeX) ---
\newcommand{\kB}{k_{\mathrm B}}
\newcommand{\NA}{N_{\mathrm A}}
\newcommand{\lamT}{\lambda_{\mathrm T}}
\newcommand{\dd}{\mathrm d}  % differential symbol
\newcommand{\ee}{\mathrm e}  % exponential base

\title{Ideal Gas Law from Statistical Thermodynamics via the 3D Particle-in-a-Box\\\large Complete, Step-by-Step Derivation}
\author{}
\date{\today}

\begin{document}
\maketitle
\begin{abstract}
We derive $PV=N\kB T$ from quantum mechanics and statistical thermodynamics without skipped steps. Starting from the 3D particle-in-a-box spectrum we build $Z_1$, pass to the continuum with controlled errors, construct $Z_N$ for indistinguishable particles, obtain $F$, and extract $P$. A classical phase-space derivation, the Sackur--Tetrode entropy, and validity conditions are included.
\end{abstract}

\section{Quantum Mechanics of a Particle in a Cubic Box}
A particle of mass $m$ confined to $V=[0,L]^3$ with infinite walls obeys
\begin{equation}
-\frac{\hbar^2}{2m}\nabla^2\psi(\bm r)=E\psi(\bm r),\qquad \psi\big|_{\partial V}=0.
\label{eq:TISE}
\end{equation}
Separation of variables gives normalized eigenfunctions and energies
\begin{equation}
E_{\bm n}=\frac{\hbar^2\pi^2}{2mL^2}\left(n_x^2+n_y^2+n_z^2\right),\qquad \bm n=(n_x,n_y,n_z)\in\mathbb N^3.
\label{eq:boxE}
\end{equation}

\section{Single-Particle Partition Function: Exact Sum and Classical Limit}
The canonical one-particle partition function is
\begin{equation}
Z_1(T,V)=\sum_{n_x=1}^{\infty}\sum_{n_y=1}^{\infty}\sum_{n_z=1}^{\infty}\exp\!\left[-\beta\,\frac{\hbar^2\pi^2}{2mL^2}(n_x^2+n_y^2+n_z^2)\right],\qquad \beta\equiv(\kB T)^{-1}.
\label{eq:Z1sumExact}
\end{equation}
Setting $a\equiv \beta\hbar^2\pi^2/(2mL^2)$, the sum factorizes:
\begin{equation}
Z_1=\Bigg[\sum_{n=1}^{\infty}\ee^{-a n^2}\Bigg]^3.
\label{eq:factor}
\end{equation}
In the classical (high-$T$/large-$L$) limit $a\ll 1$, the Riemann-sum bounds
\begin{equation}
\int_0^{\infty}\ee^{-a x^2}\,\dd x-\tfrac{1}{2}\ee^{-a}
\;<\; \sum_{n=1}^{\infty}\ee^{-a n^2}
\;<\; \int_0^{\infty}\ee^{-a x^2}\,\dd x\qquad(a>0)
\end{equation}
show the relative error vanishes as $a\to0^+$. Using
\begin{equation}
\int_0^{\infty}\ee^{-a x^2}\,\dd x=\frac{1}{2}\sqrt{\frac{\pi}{a}},
\label{eq:halfGaussian}
\end{equation}
we obtain
\begin{equation}
Z_1\simeq\left(\tfrac{1}{2}\sqrt{\tfrac{\pi}{a}}\right)^3=\frac{1}{8}\left(\frac{2mL^2}{\beta\hbar^2\pi}\right)^{3/2}.
\label{eq:Z1a}
\end{equation}
With $V=L^3$, $\beta^{-1}=\kB T$, and $h=2\pi\hbar$,
\begin{equation}
\boxed{\;Z_1=V\left(\frac{2\pi m\kB T}{h^2}\right)^{3/2}=\frac{V}{\lamT^3},\qquad \lamT\equiv\frac{h}{\sqrt{2\pi m\kB T}}\;}
\label{eq:Z1Thermal}
\end{equation}
\paragraph{Density of states route.} Counting $\bm n$-points in the first octant: $N(E)=\tfrac{1}{8}\cdot\tfrac{4}{3}\pi R^3$ with $R^2=(2mL^2/\hbar^2\pi^2)E$. Then $g(E)=\dd N/\dd E=\tfrac{V}{4\pi^2}\left(\tfrac{2m}{\hbar^2}\right)^{3/2}\!\sqrt{E}$ and $Z_1=\int_0^\infty g(E)\ee^{-\beta E}\,\dd E$ gives \cref{eq:Z1Thermal}.

\section{N Noninteracting Indistinguishable Particles}
For $N$ noninteracting particles in the Maxwell--Boltzmann regime,
\begin{equation}
\boxed{\;Z_N(T,V)=\frac{Z_1^N}{N!}=\frac{1}{N!}\left(\frac{V}{\lamT^3}\right)^N\;}
\label{eq:ZNfull}
\end{equation}
with degeneracy parameter $n\lamT^3\ll1$ and $n\equiv N/V$.

\section{Helmholtz Free Energy and Pressure}
The Helmholtz free energy is
\begin{equation}
F(T,V,N)=-\kB T\ln Z_N=-\kB T\big[N\ln V-3N\ln\lamT-\ln N!\big].
\label{eq:Fraw}
\end{equation}
Using Stirling’s expansion
\begin{equation}
\ln N!=N\ln N-N+\tfrac{1}{2}\ln(2\pi N)+\mathcal O\!\left(\frac{1}{N}\right),
\label{eq:stirling}
\end{equation}
we can write
\begin{align}
F(T,V,N)&=-N\kB T\ln\!\left(\frac{V}{N}\right)+3N\kB T\ln\lamT +\tfrac{1}{2}\kB T\ln(2\pi N)+\mathcal O\!\left(\tfrac{\kB T}{N}\right).
\label{eq:Fexpanded}
\end{align}
Since $\lamT$ is independent of $V$,
\begin{equation}
\boxed{\;P=-\left(\frac{\partial F}{\partial V}\right)_{T,N}=\frac{N\kB T}{V}\;}
\label{eq:Pressure}
\end{equation}

\section{Equation of State}
Rearranging \cref{eq:Pressure},
\begin{equation}
\boxed{\;PV=N\kB T\;}
\label{eq:IGL}
\end{equation}
and with $N=n\,\NA$, $R\equiv \NA\kB$ one gets $PV=nRT$.

\section{Classical Phase--Space Derivation}
The canonical integral for one particle is
\begin{equation}
Z_1=\frac{1}{h^3}\int_V\dd^3x\int_{\mathbb R^3}\ee^{-\beta\,\bm p^2/2m}\,\dd^3p
=\frac{V}{h^3}\left(\int_{-\infty}^{\infty}\ee^{-\beta p^2/2m}\,\dd p\right)^{\!3}.
\end{equation}
Using $\int_{-\infty}^{\infty}\ee^{-\alpha p^2}\,\dd p=\sqrt{\pi/\alpha}$ with $\alpha=\beta/(2m)$ yields $Z_1=V(2\pi m/\beta h^2)^{3/2}=V/\lamT^3$, consistent with the quantum counting.

\section{Entropy: Sackur--Tetrode}
From \cref{eq:Fexpanded},
\begin{equation}
S= -\left(\frac{\partial F}{\partial T}\right)_{V,N}
= N\kB\left[\ln\!\left(\frac{V}{N\lamT^3}\right)+\frac{5}{2}\right] + \mathcal O\!\left(\frac{1}{N}\right).
\label{eq:ST}
\end{equation}

\section{Validity}
Assumptions: noninteracting particles; classical (nondegenerate) limit $n\lamT^3\ll1$; $L\gg \lamT$ for replacing sums by integrals; and the thermodynamic limit $N,V\to\infty$ with $N/V$ fixed so that $\mathcal O(1/N)$ terms vanish.

\end{document}
